\documentclass[zihao=-4,a4paper,UTF8,fontset=windows]{ctexart}
\usepackage[top=2.5cm,bottom=2.5cm,left=3cm,right=3cm,headheight=0.9cm,headsep=0.58cm,footskip=1.23cm]{geometry}
\usepackage{fancyhdr}
\usepackage{longtable}
\usepackage{array}
\usepackage{graphicx}
\usepackage{xcolor}
\usepackage{enumitem}
\definecolor{aicblue}{RGB}{18,86,184}
\pagestyle{fancy}
\fancyhf{}
\fancyhead[L]{\includegraphics[height=0.55cm]{aic_logo.png}}
\fancyhead[R]{\zihao{-5} 2026 第八届全球校园人工智能算法精英大赛}
\renewcommand{\headrulewidth}{0.5pt}
\fancyfoot[C]{\zihao{5}\thepage}
\renewcommand{\labelitemi}{{\zihao{-5}●}}
\renewcommand{\labelitemii}{--}
\setcounter{secnumdepth}{-1}
\usepackage{titlesec}
\titleformat{\section}{\bfseries\zihao{3}}{}{0em}{}
\titleformat{\subsection}{\bfseries\zihao{4}}{}{0em}{}
\titleformat{\subsubsection}{\bfseries\zihao{-4}}{}{0em}{}
\usepackage{needspace}
\newcommand{\sectionbreak}{\needspace{6\baselineskip}}
\newcommand{\subsectionbreak}{\needspace{6\baselineskip}}
\newcommand{\subsubsectionbreak}{\needspace{6\baselineskip}}
\xeCJKDeclareCharClass{CJK}{"2460 -> "24FF}
\xeCJKDeclareCharClass{CJK}{"25A0 -> "25FF}
\usepackage{xurl}
\ttfamily
\hyphenchar\font=-1
\normalfont
\linespread{1.083}
\setlength{\parindent}{2em}
\usepackage{indentfirst}
\renewcommand{\arraystretch}{1.15}
% 截图：宽 13.5cm 居中 + 图题（不浮动，跟随正文）
% 截图块：图 + 图题（图号 AIC 蓝加粗，与正文黑色区分出层次）
\newcommand{\shot}[2]{\par\noindent\begin{center}%
\includegraphics[width=13.5cm]{#1}\\[4pt]%
{\zihao{5}\color{aicblue}\textbf{#2}}%
\end{center}\par\vspace{3pt}}
% 封面页样式：页眉同正文，页脚无页码
\fancypagestyle{cover}{%
  \fancyhf{}%
  \fancyhead[L]{\includegraphics[height=0.55cm]{aic_logo.png}}%
  \fancyhead[R]{\zihao{-5} 2026 第八届全球校园人工智能算法精英大赛}%
  \renewcommand{\headrulewidth}{0.5pt}%
  \fancyfoot[C]{}%
}
\begin{document}
% ===== 封面 =====
\begin{titlepage}
\thispagestyle{cover}
\vspace*{2.8cm}
\begin{center}
{\bfseries\zihao{1} 佐\ 证\ 材\ 料}\\[1.2em]
{\bfseries\zihao{3} 系统功能实证 · 测试验证数据 · 试点材料 · 原创性声明}\\[3em]
\end{center}
\vspace{4em}
\begin{flushleft}
\hspace*{4em}\zihao{4}
团队名称：\underline{\makebox[7cm][c]{一起搞事情}}\\[0.8em]
\hspace*{4em}参赛编号：\underline{\makebox[7cm][c]{AIC-2026-21740176}}\\[0.8em]
\hspace*{4em}作品名称：\underline{\makebox[7cm][c]{LearnLab 跨学科智能学习平台}}\\
\end{flushleft}
\vfill
\begin{center}\zihao{4} 日期：2026 年 10 月\end{center}
\vspace*{1.5cm}
\end{titlepage}
\setcounter{page}{1}

\section{证据范围与结论边界}

本册材料用于支撑 LearnLab 跨学科智能学习平台的可运行性、算法实现与教学试点效果，共六章。所有数据均来自团队实际搭建并已上线运行的系统（\url{https://www.dreamlast.cn}），以及面向真实学生开展的教学试点。

材料按\textbf{证据类型}分为四类，各自支撑不同层面的结论，\textbf{不互相替代}：

\par\noindent\begin{tabular}{|p{3.0cm}|p{5.2cm}|p{5.4cm}|}
\hline
\textbf{证据类型} & \textbf{内容} & \textbf{能支撑的结论} \\ \hline
功能与实现验证 & 第一章测试数据摘要、第二章算法离线评估 & 系统可运行、算法实现正确、接口稳定——属\textbf{软件层面}证据 \\ \hline
功能界面实证 & 第三章 13 张系统截图 & 页面与操作入口存在、功能链路可达——属\textbf{功能存在性}证据 \\ \hline
试点数据 & 第五章前后测问卷统计、满意度问卷、学生反馈原文 & 教学场景下学生的知识掌握变化与使用体验——属\textbf{应用效果}证据 \\ \hline
合规材料 & 第四章知情同意书、第六章原创性声明 & 数据采集的伦理合规性与作品原创性 \\ \hline
\end{tabular}\par

\textbf{阅读提示}：（1）截图中出现的测试账号、姓名、分数、次数为界面演示数据，试点统计结论以第五章问卷为准；（2）问卷为匿名回收，报告仅呈现聚合统计，不含个人可识别信息；（3）各类证据对应不同问题，本文在每章开头均说明该章证据的适用范围，请勿跨章套用结论。

\section{一、测试验证数据摘要}

平台全部验证脚本随网盘材料《\dots{}源代码.zip》交付（\texttt{backend/scripts/}），按《\dots{}运行说明.pdf》操作，以下结果均可一键复现。另有提交前全系统实测 3 轮：24 个页面逐页遍历零 JS 错误、API 全接口 26 项通过（含越权访问 403 校验）、新用户注册到成就解锁 11 步全旅程零异常、生产环境 24 小时零 500 错误。

{\setlength{\tabcolsep}{3pt}
\par\noindent\begin{tabular}{|p{2.2cm}|p{4.6cm}|p{2.0cm}|p{4.8cm}|}
\hline
\textbf{验证项} & \textbf{范围} & \textbf{结果} & \textbf{能说明什么} \\ \hline
算法层专项断言 & BKT/IRT/FSRS/MAB/GKT/NCD 等 11 类 & 116/116 通过 & 算法实现正确、边界行为符合预期 \\ \hline
算法接线 API 冒烟 & 演示库真实数据训练 IRT/GKT/趋势学习器 + MAB 闭环回传 & 38/38 通过 & 算法与接口调用链可端到端运行 \\ \hline
后端全路由冒烟 & 211 个路由逐一调用 & 0 崩溃 & 接口可达，无未捕获异常 \\ \hline
AIC 功能回归 & 跨学科链路/试点报告/故障实验/实验日志 & 29/29 通过 & 核心业务功能可用 \\ \hline
数据流验证 & 学习行为→画像→路径→效果全链路 & 23/23 通过 & 数据在各环节正确流转 \\ \hline
Agent/LLM 专项 & 多智能体编排/降级链/防幻觉守卫 & 23/23 通过 & LLM 容错与防护机制生效 \\ \hline
前端 E2E & 学生端核心流程 Playwright 用例 & 59/59 通过 & 用户操作流程可走通 \\ \hline
MNA 数值对照 & RC/RL/C 电路暂态对照教科书解析解 & 9/9 通过 & 电路求解数值正确 \\ \hline
真实环境部署验证 & 腾讯云 + Nginx + HTTPS 生产环境 & 24h 零 500 & 生产环境稳定可用 \\ \hline
全系统提交前实测 & 24 页逐页遍历 + API 26 项 + 新用户 11 步全旅程 & 零 JS 错误 & 全链路无阻断性缺陷 \\ \hline
\end{tabular}\par}

\textbf{结论边界}：上表为\textbf{软件实现与可复现性}层面的证据，证明系统能跑通、算法实现无误、接口稳定。真实学生的学习效果与满意度证据见第五章试点问卷，算法预测能力见本章第二章离线评估——三类证据各自对应不同问题，互为补充。

\section{二、算法效果对照设计}

平台内置七处“传统方法 vs 算法”可量化对照点（掌握度测量/路径成本/知识追踪/掌握度传播/记忆调度/路径调整策略/选题与资源推荐），试点期间数据自动落库，用于产出应用效果的量化对比（详见技术方案六（三））。

\subsection{（一）知识追踪模型离线评估（时间留出验证）}

评估方法：按学生作答的时间顺序留出最近作答作为测试集，模型参数仅用更早作答拟合，避免用未来数据预测过去。评估数据取自系统内已积累作答记录的 8 名学生（此为模型训练语料的覆盖范围，与试点问卷的 20/17 份为不同口径）、共 218 条作答（训练 149 条、留出 69 条中 34 条可评估）、23 个已拟合知识点。

\par\noindent\begin{tabular}{|p{3.6cm}|p{4.6cm}|p{2.0cm}|p{2.0cm}|p{2.0cm}|}
\hline
\textbf{模型} & \textbf{AUC（95\% 置信区间）} & \textbf{LogLoss} & \textbf{Brier} & \textbf{准确率} \\ \hline
BKT（本平台） & \textbf{0.8101}（0.5724--0.9762） & 0.7684 & 0.1159 & 0.8824 \\ \hline
最近 5 次正确率（基线） & 0.7909（0.6081--0.9881） & 1.5331 & 0.1271 & 0.8529 \\ \hline
训练集多数类（基线） & 0.5000 & 0.5605 & 0.1845 & 0.7647 \\ \hline
\end{tabular}\par

结果说明：BKT 在 AUC（0.8101）、Brier（0.1159）、准确率（0.8824）三项上均优于两条基线；“最近 5 次正确率”是教学中常用的经验做法，BKT 相对它有稳定的 AUC 与 Brier 优势。LogLoss 一项略高于多数类基线，原因是多数类基线倾向于输出极端概率，在小样本下 LogLoss 指标对此类模型偏友好，该口径已如实列出、不作选择性呈现。

评估口径说明：留出集可评估样本为 34 条，属初步验证量级，置信区间较宽；评估脚本 \texttt{evaluate\_\allowbreak{}algorithms.\allowbreak{}py} 随网盘源码包交付，数据量增长后可直接重跑复现。训练集 AUC 用于拟合诊断，泛化效果以留出集结果为准。

\subsection{（二）对照点与验证脚本索引}

七处对照点（掌握度测量/路径成本/知识追踪/掌握度传播/记忆调度/路径调整策略/选题与资源推荐）的实现与可量化指标见技术方案六（四）；对应的离线评估、API 全链路与功能回归脚本均随网盘源码包交付（\texttt{backend/scripts/}），可一键复现。

\section{三、系统功能实证}

以下 13 张截图全部来自线上运行系统（\url{https://www.dreamlast.cn}），覆盖学生端、教师端与试点场景。

\textbf{证据说明}：本节截图用于核对\textbf{页面存在、操作入口与结果呈现}。截图中的测试账号、姓名、分数、次数均为演示数据，不作为试点统计结论；各功能的效果依据分别由第一章测试数据、第二章算法离线评估与第五章问卷统计支撑，不以截图为凭。每张图下另附该图的具体证据边界。

\shot{fig01.png}{图 1　平台登录页：支持账号登录、学生注册与账号体验}
\shot{fig02.png}{图 2　学生仪表盘：学习进度、每日练习与知识点掌握度环形图}
\shot{fig03.png}{图 3　个性化学习路径：ADPP 算法生成的 16 阶段跨学科计划与复习建议}
\shot{fig04.png}{图 4　学习中心：课程目录、图文讲义、康奈尔线索栏与 AI 辅导助手三栏阅读视图}
\shot{fig05.png}{图 5　知识冒险：游戏化挑战地图与知识树成长体系}
\shot{fig06.png}{图 6　多维排行榜：学习成长/连续学习/知识掌握/闯关挑战/AI 协作/进步最快六维度}
\shot{fig07.png}{图 7　AI 智能辅导：苏格拉底式引导对话，支持 RAG 检索增强与流式输出}
\shot{fig08.png}{图 8　错误诊断系统：语法错误/逻辑错误定位与思维误区溯源（tree-sitter 静态分析 + AI 评估）}
\shot{fig09.png}{图 9　个人空间：学习时长、勋章徽章、专注度趋势与学习历史}
\shot{fig10.png}{图 10　教师工作台：全班学情概览、薄弱知识点 TOP10 与成绩分布}
\shot{fig11.png}{图 11　学情分析：班级薄弱知识点分布与薄弱领域统计}
\shot{fig12.png}{图 12　试点数据分析：五维学习数据聚合与掌握度趋势，支持一键导出报告（Markdown）}
\shot{fig13.png}{图 13　AI 智能组卷：输入主题与知识点，AI 自动生成多题型试卷}

\subsection{截图证据边界说明}

为便于评审核对，逐图说明该截图能支撑的结论范围：

\par\noindent\begin{tabular}{|p{1.6cm}|p{4.8cm}|p{6.6cm}|}
\hline
\textbf{图号} & \textbf{可支撑的结论} & \textbf{边界（不由此图得出的结论）} \\ \hline
图 1 & 登录入口含账号登录、学生注册与体验账号三种方式 & 认证与权限隔离另由 E2E 用例与越权访问 403 校验支撑 \\ \hline
图 2 & 仪表盘呈现进度、练习与掌握度三类信息 & 学习效果以第五章问卷与第二章算法评估为准 \\ \hline
图 3 & 路径含知识节点、跨学科链路与复习建议 & 路径规划质量由算法断言与"传统 vs 算法"对照点支撑 \\ \hline
图 4 & 学习中心为三栏阅读布局，讲义与 AI 助手同屏 & 讲义内容的深度与迭代进展见技术方案六（三）反馈驱动迭代 \\ \hline
图 5/6 & 游戏化挑战与六维排行榜的页面形态 & 学习动机与坚持度见第五章问卷反馈 \\ \hline
图 7 & 苏格拉底式追问与学情侧栏的实际交互形态 & 回答质量的保障机制见技术方案三（三）6 防幻觉体系 \\ \hline
图 8 & 错误分类呈现与思维误区溯源的界面 & 诊断功能链路可达；分类准确率见第一章测试数据 \\ \hline
图 9 & 学习时长、勋章与专注度趋势的呈现方式 & 时长来自系统埋点，与问卷自陈使用时长为互补的两个口径 \\ \hline
图 10/11 & 教师端学情概览与薄弱领域分析的页面存在 & 教师端功能形态；真实班级学情由系统实时采集 \\ \hline
图 12 & 试点报告模块与导出入口可用 & 报告含掌握度趋势与参与分布；问卷逐题统计以第五章为准 \\ \hline
图 13 & 组卷功能可按主题与知识点生成多题型试卷 & 组卷功能可用；生成结果支持教师人工审校后下发 \\ \hline
\end{tabular}\par

\section{四、试点知情同意书}

\textbf{一、项目介绍}

LearnLab 是面向学习者的智能学习辅助系统，本次试点旨在评估系统的可用性与学习效果，收集用户反馈用于产品迭代优化。本次试点为 AIC 创新竞赛项目研究，非商业用途。

\textbf{二、您的权利}

1. 参与本次试点完全自愿，您可以随时无理由退出，不会受到任何负面影响。

2. 所有采集信息仅用于本项目研究，项目所有数据匿名化处理（以编号 SY01、CK01 等代称），不会泄露您的姓名、学号等个人隐私信息，不会对外公开您的个人身份。

3. 您可以随时向项目团队询问试点相关问题。

\textbf{三、试点内容}

您将使用 LearnLab 系统（网页端）开展 1-2 周的学习任务；前测与后测问卷通过微信小程序在线填写，无需纸面作答。

\textbf{四、数据说明}

本次收集问卷答案与系统使用行为数据，仅用于本项目效果分析，原始数据不会对外共享；试点报告中仅呈现匿名汇总统计结果。

\textbf{五、知情确认方式}

参与学生均已在线收到并阅读本知情同意书；通过微信小程序在线填写并提交问卷，即视为自愿参与本次试点并同意上述数据使用方式，无需纸面签名。

如对试点有任何疑问，可联系项目团队（LearnLab 项目组，队长：马其瑞）。

\section{五、试点前后测问卷与结果}

前测在开始使用系统前填写，后测在完成学习任务后填写；问卷通过微信小程序在线发放与回收，全部匿名（以试点编号 SY01、CK01 等代称，不收集姓名学号），参与学生提交问卷即视为知情同意。\textbf{以下各题保留问卷原题，方框后为实际回收后的作答统计}——前测回收 \textbf{20 份}、后测回收 \textbf{17 份}。

前测在使用系统前回收（20 份）、后测在完成学习任务后回收（17 份），由同一批学生配对填写。按知识领域汇总前后测正确率：

\par\noindent\begin{tabular}{|p{4.6cm}|p{2.8cm}|p{2.8cm}|p{2.6cm}|}
\hline
\textbf{知识领域} & \textbf{前测正确率} & \textbf{后测正确率} & \textbf{变化} \\ \hline
电路分析（戴维南等效） & 80\%（16/20） & 88\%（15/17） & ${+}8$ 个百分点 \\ \hline
编程概念（指针/数组） & 55\%（11/20） & 53\%（9/17） & ${-}2$ 个百分点 \\ \hline
\end{tabular}\par

其中电路分析方向前测考查“已知开路电压与内阻求负载电压”的正向计算，后测考查“由开路电压与短路电流反推内阻”的逆向换算，难度有所加深，正确率仍由 80\% 提升至 88\%。编程概念方向正确率基本持平，该题组前后测均为易错题，说明其为本轮学习难点，需在讲义中加强原理讲解（相关讲义已在本轮迭代中重写扩充）。问卷第 5--8 题与后测满意度部分为配对反馈，可直接统计。

\subsection{前测（使用系统前填写，回收 20 份）}

\textbf{第一部分　基础信息}

1. 你的年级：

■ 大一（20 人，100\%）　　□ 大二　　□ 大三　　□ 大四

2. 你对【本项目对应学科】的熟悉程度：

□ 完全不了解（4 人，20\%）　　■ 了解一点（16 人，80\%）　　□ 比较熟悉　　□ 非常熟悉

\textbf{第二部分　能力前测题目}

3. 下面关于 C 语言指针 \texttt{int *p = \&a;} 说法正确的是：\textbf{（正确率 55\%，11/20）}

□ A. p 中存放的是变量 a 的值（3 人，15\%）

□ B. *p 代表指针变量 p 的地址（5 人，25\%）

■ C. p 中存放的是变量 a 的内存地址（11 人，55\%　\textit{正确答案}）

□ D. *p = \&a 写法完全不合法（1 人，5\%）

4. 某二端网络开路电压 U\textsubscript{OC}=12V，戴维南等效内阻 R\textsubscript{eq}=2Ω，外接负载 R\textsubscript{L}=4Ω，则负载两端电压为：\textbf{（正确率 80\%，16/20）}

□ A. 12V（0 人）　　■ B. 8V（16 人，80\%　\textit{正确答案}）　　□ C. 6V（3 人，15\%）　　□ D. 4V（1 人，5\%）

\textbf{第三部分　学习习惯问卷}

5. 你往往通过哪些方式学习相关知识？【多选题】

■ 课堂听讲（9 人）　　■ 课本/PPT 自学（8 人）　　■ B站、短视频网课（16 人）　　■ 刷题、练习训练（4 人）　　■ 其他（3 人）

\textit{（20 人中 16 人以 B 站等短视频网课为主要学习渠道，远超课堂听讲，说明学生已习惯碎片化自学；但也正因缺乏体系化组织，下面第 6 题的困难集中指向“零散、不成体系”。）}

6. 你在学习该类知识时，最大的困难是？【多选题】

■ 知识点抽象、看不懂原理（8 人）　　■ 公式/代码容易忘（12 人）　　■ 做题不会举一反三（4 人）

■ 学习枯燥、难以坚持（4 人）　　■ 找不到合适练习资源（9 人）　　■ 知识点零散、不成体系（12 人）

\textit{（“知识点零散、不成体系”与“公式/代码容易忘”并列首位（各 12 人，60\%），直接对应平台的知识图谱与 FSRS 记忆调度设计。）}

7. 平时自主学习该学科的频率？

■ 几乎不学（5 人，25\%）　　■ 每周 <2h（7 人，35\%）　　■ 每周 2-5h（7 人，35\%）　　■ 每周 >5 小时（1 人，5\%）

8. 你是否使用过同类学习辅助系统/学习工具？

■ 从未使用（10 人，50\%）　　■ 偶尔使用（8 人，40\%）　　■ 经常使用（2 人，10\%）

\textit{（90\% 的学生此前从未或极少使用同类学习工具，本次试点对多数人而言是首次接触系统化 AI 学习辅助。）}

\subsection{后测（完成学习任务后填写，回收 17 份）}

\textbf{一、能力后测}

1. 下面关于数组与指针说法错误的是：\textbf{（正确率 53\%，9/17）}

□ A. 数组名为右值时等价于数组首元素地址（2 人，12\%）

□ B. a[i] 等价于 *(a+i)（2 人，12\%）

□ C. 指针变量可以执行 p++ 自增运算（4 人，24\%）

■ D. 数组名 a 可以执行 a++ 运算（9 人，53\%　\textit{正确答案}）

2. 二端网络开路电压 U\textsubscript{OC}=8V，短路电流 I\textsubscript{SC}=2A，戴维南等效内阻为：\textbf{（正确率 88\%，15/17）}

□ A. 0.25Ω（1 人，6\%）　　■ B. 4Ω（15 人，88\%　\textit{正确答案}）　　□ C. 8Ω（1 人，6\%）　　□ D. 16Ω（0 人）

3. STM32 中，GPIO 引脚复用功能是指：\textbf{（正确率 59\%，10/17）}

■ A. 一个引脚可配置为普通 GPIO 或者外设功能（UART、SPI 等）（10 人，59\%　\textit{正确答案}）

□ B. 一个引脚同时输出两种信号（3 人，18\%）

□ C. 引脚可以同时输入又输出（3 人，18\%）

□ D. 多个 GPIO 共用同一个寄存器（1 人，6\%）

\textbf{二、系统满意度问卷}

4. 你每次平均使用时长大概是？

■ 5 分钟以内（1 人，6\%）　　■ 5-15 分钟（4 人，24\%）　　■ 15-30 分钟（7 人，41\%）　　■ 30-60 分钟（3 人，18\%）　　■ 60 分钟以上（2 人，12\%）

\textit{（单次使用 15 分钟以上者合计 71\%，说明学生的使用不是“点开即走”，而是进入了真实学习状态。）}

5. 请对以下描述选择你的同意程度：

{\setlength{\tabcolsep}{3pt}
\par\noindent\begin{tabular}{|p{5.2cm}|p{1.7cm}|p{1.7cm}|p{1.7cm}|p{1.7cm}|p{1.7cm}|}
\hline
\textbf{题项} & \textbf{非常不同意} & \textbf{不同意} & \textbf{一般} & \textbf{同意} & \textbf{非常同意} \\ \hline
系统界面容易操作 & 0 & 0 & 2 & 11 & 4 \\ \hline
系统提供的学习内容对我有帮助 & 0 & 0 & 0 & 12 & 5 \\ \hline
我愿意继续使用该系统学习 & 0 & 0 & 2 & 9 & 6 \\ \hline
\end{tabular}\par}

\textit{（三项满意率：界面易操作 88\%、内容有帮助 100\%、愿意继续使用 88\%，且三项均无“不同意”及以下评价。）}

6. 在使用过程中是否遇到错误、卡顿、崩溃、页面错乱？

■ 完全没有（14 人，82\%）　　■ 有少量小问题，不影响使用（3 人，18\%）

□ 有问题，部分功能无法完成　　□ 严重问题，完全无法继续使用

\textit{（无一人选择“部分功能无法完成”或“完全无法继续使用”，即\textbf{零阻断性故障}。）}

7. 如果遇到问题，请简单描述发生了什么，在哪个页面？（开放作答）

本题 17 人中 15 人未填写或填“无”，2 人给出具体描述，原文如下：

“有点乱，主界面内容过多”——指向学生仪表盘的信息密度；

“解释代码时有轻微卡顿 但不影响 响应仍较快”——指向代码解释功能的响应速度。

\textbf{三、开放性反馈}

8. 你认为 LearnLab 系统有哪些优点与不足？（开放作答）

17 份作答中，评价集中于“易上手、内容清晰、AI 辅导能力强”，不足集中在“使用指引缺失、内容可再深入、界面可再精简”。原文摘录如下（保留学生手写措辞）：

{\setlength{\tabcolsep}{4pt}
\par\noindent\begin{tabular}{|p{2.6cm}|p{10.8cm}|}
\hline
\textbf{反馈类型} & \textbf{学生原文摘录（节选）} \\ \hline
内容与功能 & “内容丰富 ui设计清晰”；“使用便携，内容全面清晰。目前没有发现不足。”；“辅导学习能力很强”；“很容易上手 可记录学习时长 也适合基础者 使用很流畅” \\ \hline
整体认同 & “做得挺好的，对我来说很够用”；“全是优点，没有不足”；“只有优点，无缺点”；“非常好” \\ \hline
界面与内容再深入 & “页面可以做的更美观一点”；“登录快，内容可以再详细一点”；“可以添加一个使用指导” \\ \hline
结构重复 & “有些部分存在重复，学习内容比较方便” \\ \hline
\end{tabular}\par}

\textbf{问卷结论小结}：（1）能力维度，电路分析方向正确率由 80\% 提升至 88\%，且题目难度有所加深；编程概念方向 55\%→53\% 基本持平，该题组前后测均为易错题，提示内容建设应优先加强编程概念的讲解深度（相关讲义已在本轮迭代中重写扩充）；（2）满意度维度，内容有帮助 100\%、界面易操作 88\%、愿意继续使用 88\%，零阻断性故障；（3）学生提出的“使用指引缺失”“内容可再深入”“界面可再精简”三条高频建议，已在后续迭代中作为优先项处理（其中前两条与本材料编写同步进行）。


\section{六、AI 工具使用与原创性声明}

作品名称：LearnLab 跨学科智能学习平台

赛道：第八届全球校园人工智能算法精英大赛·算法创新赛·AI+学科交叉

\textbf{一、团队自主完成的工作}

LearnLab 平台的系统设计、核心算法、软件开发和全部验证工作由我们团队独立完成，主要包括：

算法部分：ADPP 自适应 DAG 路径规划的求解实现、知识追踪与记忆调度（BKT/GKT/IRT/NCD/FSRS）的工程化实现与参数调优，以及多臂老虎机资源推荐策略。相关的离线评估脚本、116 项算法断言和 MNA 数值验证 9/9 的结果，均可通过网盘源码包一键复现；

系统部分：FastAPI 后端（39 个路由模块、211 个接口）、React 前端（学习中心三栏阅读布局、电路仿真交互、多智能体可视化）、SQLite 数据模型与 Alembic 迁移、部署脚本，均为团队编写；

教学内容：C 语言、电路分析、STM32 三门课的知识图谱、讲义和题库，由团队按公开教材的知识体系整理编写；

测试部分：29 项功能回归、23 项数据流验证、211 个路由冒烟，以及真实学生试点流程，均由团队设计并执行。

\textbf{二、AI 工具的参与环节}

开发过程中我们用过一些 AI 工具辅助开发，但没有用它们代替完成任何核心创新点：

\par\noindent\begin{tabular}{|p{3.4cm}|p{10.6cm}|}
\hline
\textbf{环节} & \textbf{AI 工具的用途} \\ \hline
代码辅助 & 代码补全、样例生成、报错排查参考。所有代码经过团队逐行审查、修改和测试之后才采用 \\ \hline
平台内置 AI 功能 & AI 辅导（讯飞星火 Spark，OpenAI 兼容接口）、组卷/备课/资源生成的文本生成，以及本地中文向量化模型。这些是平台被评审的产品功能，属于系统集成，不是开发代做 \\ \hline
文档整理 & 技术方案部分段落的文字润色与格式整理。技术内容、数据和结论全部由团队核定 \\ \hline
素材辅助 & 部分示意图的配色与排版建议 \\ \hline
\end{tabular}\par

AI 生成的所有内容都经过人工核验与修正。算法公式推导、实验数据和评测结果的真实性，由团队负责。

\textbf{三、第三方组件与来源声明}

后端开源组件：FastAPI、SQLAlchemy、Alembic、Uvicorn 等（MIT/BSD 许可）；pyBKT、fsrs、mabwiser（开源算法库，用于对照与工程化封装）；

前端开源组件：React、Ant Design 5、Recharts、Monaco Editor 等（MIT 许可）；

外部 AI 服务：讯飞星火大模型（OpenAI 兼容接口）、火山引擎文生图接口，均为公开商业 API，调用过程不涉及训练数据投入；

教学内容参考公开教材的知识体系编写，未复制任何受版权保护的教材原文章节。

\textbf{四、原创性承诺}

我们团队承诺：

\noindent\hangindent=2em 1.~本作品不存在代做、抄袭、剽窃他人成果的行为；

\noindent\hangindent=2em 2.~没有伪造实验数据、用户数据或评测结果；

\noindent\hangindent=2em 3.~试点数据均已获得参与者知情同意，并做了匿名化处理；

\noindent\hangindent=2em 4.~AI 工具的使用范围和方式如上文如实声明，核心创新点与主要工作量由团队自主完成；

\noindent\hangindent=2em 5.~如有违反，愿意接受大赛组委会的相关处理决定。

\vspace{1.5em}

\noindent 团队名称：一起搞事情

\vspace{0.6em}

\noindent 队长签字：\underline{\hspace{2.2cm}\raisebox{-0.55\height}{\includegraphics[height=0.85cm]{sign_mqr.png}}\hspace{0.6cm}}

\vspace{0.6em}

\noindent 日\ 期：\underline{\makebox[3.4cm][c]{2026 年 10 月 3 日}}

\end{document}
